% called by main.tex
%
\chapter{Estado de la cuestión}
\label{ch::estado}

Este capítulo sitúa el trabajo en la literatura. Se organiza en las áreas que confluyen
en el problema —mercados de predicción y el estudio académico reciente de Polymarket,
aprendizaje profundo sobre libros de órdenes, modelos fundacionales y auto-supervisados
de series temporales, microestructura y ejecución, validación temporal en finanzas, y
cambio de distribución— y termina señalando el hueco concreto que este trabajo pretende
cubrir.

\section{Trabajos previos y estudios relevantes}

Esta sección recorre la literatura relevante para el problema, organizada en bloques
temáticos. En cada bloque se resume qué proponen los trabajos de referencia y qué se toma
de ellos para este trabajo. La predicción del movimiento de precios a partir de la
microestructura es un campo activo, con propuestas asentadas tanto para mercados
tradicionales como para criptoactivos; el interés aquí es seleccionar de esa literatura los
elementos aplicables a un mercado de predicción sobre Bitcoin.

\subsection{Mercados de predicción}

Los mercados de predicción se han estudiado, sobre todo, como mecanismos de agregación de
información dispersa. Wolfers y Zitzewitz~\cite{wolfers2004prediction} muestran que pueden
producir previsiones notablemente precisas, si bien su interpretación depende del diseño
del contrato, de la liquidez y de los incentivos de los participantes. La mayor parte de
esta literatura analiza la calibración del precio frente a la resolución final del evento.
El presente trabajo se sitúa un escalón por debajo, en el terreno de la microestructura:
la dinámica intradía del precio del contrato a escala de segundos, donde aparecen
\emph{spreads}, profundidad limitada y desajustes transitorios respecto a las referencias
externas. Conviene además precisar el mecanismo: a diferencia de los mercados de predicción
basados en creadores de mercado automáticos (que fijan el precio mediante una fórmula),
Polymarket emplea un libro de órdenes, lo que hace aplicable el instrumental clásico de la
microestructura y justifica el enfoque de este trabajo.

\subsection{Polymarket como objeto de estudio académico}

Hasta 2025, Polymarket apenas contaba con literatura académica propia; en la primera
mitad de 2026 han aparecido cinco estudios que lo convierten en objeto de análisis
directo, y que delimitan con precisión el hueco que ocupa este trabajo.
% TODO M2: figura de línea temporal de esta literatura (feb-jun 2026).

Dubach~\cite{dubach2026anatomy} construye un archivo \emph{tick-level} del canal
público del libro de órdenes (30~mil millones de eventos a lo largo de 52 días) y
extrae de él ocho hechos estilizados de la microestructura de Polymarket:
concentración de la actividad en pocos mercados, \emph{spreads} dependientes del
nivel de precio, actividad de cancelación dominante y ciclos intradía marcados,
entre otros. Es el trabajo más cercano al presente en objeto y escala temporal, y
sirve aquí como referencia de contraste para el análisis exploratorio del
Capítulo~\ref{ch::metodos}: los hechos estilizados que aquel archivo documenta
sobre el conjunto de la plataforma se comparan con lo observado en el corpus propio
de mercados horarios de Bitcoin. Tsang y Yang~\cite{tsang2026election} analizan el
mercado de la elección presidencial estadounidense de 2024 con datos de la cadena
Polygon y un marco contable a nivel de transacción, incluyendo la detección de
\emph{wash trading}; su interés es forense y agregado, no operativo. Cheng, Yang y
Zou~\cite{cheng2026arbitrage} buscan oportunidades de arbitraje en los mercados de
baloncesto y encuentran que las anomalías de mercado único son rarísimas y de vida
mediana de pocos segundos — una advertencia empírica directa contra las estrategias
\emph{taker} de latencia que este trabajo confirma por otra vía.
Nechepurenko~\cite{nechepurenko2026nonretail} estudia la composición del flujo
ejecutado y muestra que la liquidez de Polymarket descansa en niveles de
participantes no minoristas con firmas de microestructura diferenciadas; para este
trabajo, ese resultado describe al adversario contra el que compite cualquier
política \emph{maker}. Por último, Qin y Yang~\cite{qin2026polymarketv1} publican la
base Polymarket-v1, el archivo on-chain completo de la primera generación del
protocolo (1,2~mil millones de operaciones entre 2022 y 2026), orientada a la
historia larga de operaciones liquidadas.

La posición de este trabajo frente a esa literatura se resume en dos diferencias.
La primera es el dato: los archivos citados capturan operaciones on-chain o el canal
público de eventos, mientras que el corpus propio sincroniza el libro L2 completo a
malla fija de 2 segundos con referencias externas simultáneas (\emph{spot},
perpetuos y oráculo) y con un contrato de latencia medido sobre el ciclo real de
órdenes — es decir, contiene exactamente la información que hace falta para simular
ejecución, no solo para describir actividad. La segunda es la pregunta: ninguno de
los cinco estudios evalúa económicamente una política de negociación congelada bajo
coste, latencia y régimen de comisiones medidos, ni valida prospectivamente sus
resultados; ese es el terreno de este trabajo.

\subsection{Aprendizaje profundo sobre libros de órdenes}

La predicción del movimiento de precios a partir del libro de órdenes es un problema con
una literatura ya madura. El trabajo de referencia es DeepLOB~\cite{zhang2019deeplob}, que
representa el libro como una matriz de dos dimensiones (niveles de precio frente a tiempo) y
aplica primero módulos convolucionales —que capturan la estructura espacial entre los
niveles de compra y de venta— y, a continuación, capas recurrentes que modelan la
dependencia temporal; se valida sobre FI-2010~\cite{ntakaris2018}, un conjunto de referencia
público de libros de órdenes ampliamente utilizado como banco de pruebas. En la misma línea,
Tsantekidis et al.~\cite{tsantekidis2017} comparan arquitecturas convolucionales y
recurrentes para detectar cambios de precio en el libro. Sirignano y
Cont~\cite{sirignano2019universal} aportan evidencia de la existencia de características
universales en la formación de precios que el aprendizaje profundo es capaz de capturar, y
Jha et al.~\cite{jha2020digital} trasladan estas ideas al terreno de los criptoactivos,
mostrando que el enfoque no es exclusivo de los mercados tradicionales. La línea sigue
activa: TLOB~\cite{berti2025tlob} sustituye el par convolucional-recurrente por un
transformador con doble atención (espacial sobre los niveles del libro y temporal sobre
la secuencia) y documenta, además, que el rendimiento de todos estos modelos se degrada
al acercarse a condiciones de mercado realistas — observación que anticipa el resultado
central de este trabajo.

De esta literatura este trabajo adopta una decisión de diseño —tratar el libro de órdenes
como un objeto con estructura espacio-temporal, y no como una tabla plana— y, sobre ella,
introduce dos cambios. El primero es metodológico: evaluar cada modelo bajo un modelo de
ejecución realista (coste y latencia), y no solo con métricas de clasificación como es
habitual en estos trabajos. El segundo es de dominio: aplicar el enfoque a un mercado de
predicción sobre Bitcoin —y no a un mercado de acciones—, lo que obliga a combinar el libro
con referencias externas. Esa doble diferencia es, precisamente, donde este trabajo pone el
acento.

\subsection{Modelos fundacionales de series temporales y aprendizaje auto-supervisado}

Una segunda corriente reciente propone trasladar a las series temporales la receta de
los grandes modelos de lenguaje: pre-entrenar un modelo genérico sobre colecciones
masivas y heterogéneas de series y aplicarlo después a tareas nuevas sin entrenamiento
(\emph{zero-shot}) o con una adaptación ligera (\emph{fine-tuning}).
MOMENT~\cite{goswami2024moment} es el representante que este trabajo adopta como
peldaño de su escalera de modelos: una familia abierta de transformadores
pre-entrenados por reconstrucción enmascarada, con variantes que caben en hardware
modesto. En la misma familia, Chronos~\cite{ansari2024chronos} discretiza los valores
de la serie y los trata literalmente como \emph{tokens} de un modelo de lenguaje;
TimesFM~\cite{das2024timesfm} es el fundacional de predicción de Google entrenado
sobre cientos de miles de millones de puntos reales, y la línea
Moirai~\cite{liu2025moirai2} —ya en su segunda versión, con predicción cuantílica—
representa el estado del arte de los fundacionales universales de predicción. Fuera de
las series, TabPFN~\cite{hollmann2025tabpfn} lleva la misma idea al dato tabular:
una red pre-entrenada sobre el prior de millones de conjuntos sintéticos que
clasifica conjuntos pequeños por inferencia en contexto, sin entrenar. Su inclusión
en la escalera de este trabajo permite formular la pregunta simétrica —¿qué aporta
un fundacional \emph{del tipo de dato que ganó} (tabular), y no solo del tipo de
dato aparente (serie temporal)?— cuyos resultados se presentan en el
Capítulo~\ref{ch::resultados}.

El pre-entrenamiento auto-supervisado es el mecanismo que sostiene a estos modelos y
una línea propia de literatura. Las estrategias dominantes son la reconstrucción de
tramos enmascarados —SimMTM~\cite{dong2023simmtm} la formula como suma ponderada de
series vecinas fuera de la variedad de la serie original—, el contraste jerárquico
entre vistas aumentadas —TS2Vec~\cite{yue2022ts2vec}, generalizado por
SoftCLT~\cite{lee2024softclt} con asignaciones blandas que respetan la cercanía
temporal— y los objetivos híbridos: TimeSiam~\cite{dong2024timesiam} empareja
subseries pasadas y presentes con codificadores siameses, y
TimeDART~\cite{wang2025timedart} combina difusión y auto-regresión para capturar a la
vez la dinámica global y los patrones locales. Para este trabajo, esta literatura
cumple dos papeles: motiva el peldaño auto-supervisado de la escalera (un codificador
convolucional-temporal pre-entrenado por reconstrucción sobre el propio corpus) y
aporta la advertencia, repetida en sus propias evaluaciones, de que la ventaja del
pre-entrenamiento depende críticamente del volumen y de la afinidad de dominio de los
datos — condición que un mercado joven, con semanas de historia, no cumple de partida.
% TODO M2: figura-mapa de la familia de modelos (fundacionales vs SSL vs especialista).

\subsection{Microestructura, costes de ejecución y latencia}

La distancia entre predecir bien y obtener beneficio es uno de los temas clásicos de la
microestructura de mercado. Kearns y Nevmyvaka~\cite{kearns2013} revisan el papel del
aprendizaje automático en microestructura y negociación de alta frecuencia, y subrayan que
la ejecución, el coste y la latencia condicionan de forma decisiva la viabilidad de una
señal. Este trabajo cuantifica ese fenómeno en el contexto concreto de Polymarket: se mide
la latencia real del ciclo de orden y se demuestra que una señal direccional fuerte se
vuelve no rentable bajo una latencia de entrada modesta. El cambio que se introduce respecto
a esta literatura es, de nuevo, metodológico: en lugar de asumir una latencia o un coste
fijos, ambos se determinan a partir de los datos —la latencia, de forma empírica sobre el
ciclo de orden real; el coste, de forma relativa al estado del libro en cada instante—, de
modo que la evaluación económica refleje las condiciones efectivas de ejecución.

\subsection{Aprendizaje automático en finanzas y validación temporal}

El aprendizaje automático aplicado a la predicción financiera es especialmente propenso a
resultados ilusorios cuando la evaluación no respeta la estructura temporal de los datos.
Gu, Kelly y Xiu~\cite{gu2020empirical} ofrecen un marco riguroso de comparación de modelos
con validación fuera de muestra. López de Prado~\cite{lopezdeprado2018} documenta los
mecanismos por los que se cuela información del futuro y por los que se sobreajustan los
\emph{backtests}, y propone protocolos de validación que respetan la causalidad temporal.
Bailey y López de Prado~\cite{bailey2014deflated} advierten de un riesgo adicional: si se
prueban muchos modelos y se elige el mejor, las métricas de rendimiento salen infladas casi
por construcción. Estas referencias motivan tres decisiones metodológicas de este trabajo:
la partición estrictamente temporal, el protocolo fuera de pliegue (\emph{out-of-fold}) y
la congelación del candidato antes de observar el conjunto de validación fuera de muestra.

\subsection{Cambio de distribución y adaptación de dominio}

Los mercados financieros no son estacionarios, por lo que la distribución de los datos
cambia entre el periodo de entrenamiento y el de despliegue. Para amortiguar ese efecto
existen técnicas de adaptación de dominio, como la alineación de correlaciones
CORAL~\cite{sun2016coral}. En el Capítulo~\ref{ch::resultados} se mostrará un resultado
contraintuitivo: con un horizonte de datos reducido, estas técnicas no superan a una
estrategia mucho más simple, consistente en restringir la operativa al régimen en el que
el modelo fue entrenado. Además de CORAL, en este trabajo se evalúan otras estrategias de
adaptación habituales —filtros conformales, \emph{ensembles} y reentrenamiento incremental
con ventana móvil—, cuyos resultados se detallan en el Capítulo~\ref{ch::resultados}. Por
último, en cuanto a arquitecturas profundas para series temporales multivariantes, trabajos
como el \emph{Temporal Fusion Transformer}~\cite{lim2021tft} marcan una línea de evolución
natural para cuando se disponga de muchos más datos.

\subsection{Rigor de validación: pre-registro y control del sesgo de selección}

Las referencias de validación temporal citadas más arriba comparten un diagnóstico: el
riesgo dominante en la investigación financiera empírica no es el error del modelo,
sino el del investigador que prueba muchas variantes y reporta la mejor. Harvey, Liu y
Zhu~\cite{harvey2016cross} cuantifican el problema en la literatura de factores —la
mayoría de los resultados publicados no sobreviviría umbrales de contraste múltiple— y
Bailey y López de Prado~\cite{bailey2014deflated} ofrecen la corrección
correspondiente para estrategias de negociación, la ratio de Sharpe deflactada, que
penaliza explícitamente el número de ensayos. Ambas líneas atacan el sesgo
\emph{a posteriori}, corrigiendo la métrica una vez cometido el pecado.

La respuesta \emph{a priori} es el pre-registro: comprometer hipótesis, protocolo y
criterios de éxito antes de observar los datos de evaluación, práctica hoy estándar en
las ciencias experimentales~\cite{nosek2018preregistration}. Su punto débil clásico es
la verificabilidad — nada impide redactar el «pre»-registro después. Este trabajo lo
resuelve con un sello criptográfico de tiempo: el documento de pre-registro y la
política congelada se atestan con OpenTimestamps~\cite{opentimestamps} sobre la cadena
de Bitcoin, de modo que cualquier lector puede verificar que el compromiso es anterior
a los datos de la ventana de evaluación, sin confiar en el autor ni en un tercero. La
combinación —candidato congelado con hashes verificables, ventana de evaluación
prospectiva intocada hasta un corte fijado de antemano, y reporte de todos los brazos
comprometidos, favorables o no— es, hasta donde alcanza la revisión de este capítulo,
inédita en la literatura de mercados de predicción, y constituye la aportación
metodológica central del trabajo (el protocolo se detalla en el
Capítulo~\ref{ch::metodos}).

\section{Limitaciones en la literatura}

Del repaso anterior se desprenden tres limitaciones recurrentes que enmarcan la aportación
de este trabajo. En primer lugar, la evaluación predominante es de \emph{clasificación}
(precisión, AUC), y rara vez se traduce a un resultado económico bajo un modelo de
ejecución realista; como se verá, esa traducción cambia las conclusiones por completo. En
segundo lugar, gran parte de los resultados positivos se obtienen sin una latencia de
entrada explícita, cuando la latencia es precisamente lo que destruye muchas señales de
corto plazo. En tercer lugar, los entornos de criptoactivos jóvenes ofrecen pocos datos,
lo que agrava el riesgo de sobreajuste y de resultados ilusorios si la validación no es
estrictamente temporal.

El hueco que cubre este trabajo es, por tanto, metodológico: evaluar la predicción de
corto plazo \emph{bajo coste y latencia medida}, con validación estrictamente temporal y
congelación previa del candidato, y reportar cada resultado en ticks netos acompañado de
su soporte y su incertidumbre.
